\question
Один маленький диктатор решил провести реформу языка в своей стране. По его задумке в алфавите остаются только буквы $A$, $B$, $C$, $D$; причём если в последовательности есть хоть одна буква $A$, то тогда в ней обязательно должны встретиться две буквы $A$ подряд, иначе эта последовательность -- не слово.

Парочку примеров он всё же потрудился написать:

CDCD, AAABAC, DDAA -- слова.

DADA, DDDA -- не слова.

Вам было поручено найти количество слов длины 8 после реформы.

---------------

Автор -- Константин Васильев
